\section*{Hoofdstuk \ref{ch:g12:deriv} --- Afgeleiden en convexiteit}

\begin{solution}{exo:g12:deriv:1}
$f'(x) = 3x^2 - 5$.

Quotiëntregel:
$g'(x) = \dfrac{(x^2+1) - x\cdot 2x}{(x^2+1)^2} = \dfrac{1 - x^2}{(x^2+1)^2}$.

Kettingregel met $u = 2x+1$: $h'(x) = 5 \cdot 2 \cdot (2x+1)^4 =
10(2x+1)^4$.

Kettingregel met $u = x^2 + 1$:
$k'(x) = \dfrac{2x}{2\sqrt{x^2+1}} = \dfrac{x}{\sqrt{x^2+1}}$.
\end{solution}

\begin{solution}{exo:g12:deriv:2}
$f'(x) = 2x - 3$, dus $f(2) = -1$ en $f'(2) = 1$: de \omterm{def:g12:deriv:tangent}{raaklijn} is
$y = -1 + (x - 2) = x - 3$. De kloof is
\[
f(x) - (x - 3) = x^2 - 4x + 4 = (x-2)^2 \geq 0 ,
\]
dus ligt de kromme overal boven de \omterm{def:g12:deriv:tangent}{raaklijn} en raakt ze die alleen in
$x = 2$ (zoals verwacht: $f$ is \omterm{def:g12:deriv:convex}{convex}).
\end{solution}

\begin{solution}{exo:g12:deriv:3}
Allebei zijn het differentiequotiënten. Met $f(x) = (1+x)^{100}$ in $0$:
$f'(x) = 100(1+x)^{99}$, dus is de limiet $f'(0) = 100$. Met
$g(x) = \sqrt{x}$ in $4$: $g'(x) = \frac{1}{2\sqrt{x}}$, dus is de limiet
$g'(4) = \frac14$.
\end{solution}

\begin{solution}{exo:g12:deriv:4}
\omterm{def:g11:func:function}{Domein} $\R \setminus \{1\}$. Limieten: in $\pm\infty$ gedraagt $f(x)$ zich
als $x \to \pm\infty$; in $1^\pm$ streeft de teller naar $4 > 0$ en de noemer
naar $0^\pm$, dus $f \to \pm\infty$; de rechte $x = 1$ is een \omterm{def:g12:limcont:asymptote}{verticale
asymptoot}.

\[
f'(x) = \frac{2x(x-1) - (x^2+3)}{(x-1)^2} = \frac{x^2 - 2x - 3}{(x-1)^2}
= \frac{(x-3)(x+1)}{(x-1)^2}.
\]
Bijgevolg is $f' > 0$ op $\intoo{-\infty}{-1}$ en $\intoo{3}{+\infty}$, en
$f' < 0$ op $\intoo{-1}{1}$ en $\intoo{1}{3}$. De \omterm{def:g11:func:function}{functie} stijgt tot een
\omterm{def:g12:deriv:extremum}{lokaal maximum} $f(-1) = -2$, daalt naar $-\infty$, springt na de \omterm{def:g12:limcont:asymptote}{asymptoot}
naar $+\infty$, daalt tot een lokaal \omterm{def:g10:functions:extrema}{minimum} $f(3) = 6$, en stijgt daarna.
\end{solution}

\begin{solution}{exo:g12:deriv:5}
Voor $x \in \intoo{0}{15}$ heeft de doos een vierkant grondvlak met zijde
$30 - 2x$ en hoogte $x$, dus
\[
V(x) = x(30 - 2x)^2 .
\]
$V'(x) = (30-2x)^2 + x \cdot 2(30-2x)(-2) = (30-2x)(30 - 2x - 4x)
= (30-2x)(30-6x)$. Op $\intoo{0}{15}$ is $30 - 2x > 0$, dus heeft $V'$ het
teken van $30 - 6x$: positief vóór $x = 5$, negatief erna. Het volume is
maximaal voor $x = 5$ cm, met $V(5) = 5 \times 20^2 = 2000\ \text{cm}^3$.
\end{solution}

\begin{solution}{exo:g12:deriv:6}
\emph{1.} $f'(x) = 4x^3 - 12x^2$ en $f''(x) = 12x^2 - 24x = 12x(x-2)$. Dus is
$f'' > 0$ op $\intoo{-\infty}{0}$ en op $\intoo{2}{+\infty}$ (\omterm{def:g12:deriv:convex}{convex}), en
$f'' < 0$ op $\intoo{0}{2}$ (\omterm{def:g12:deriv:convex}{concaaf}). $f''$ verandert van teken in $0$ en in
$2$: allebei \omterm{def:g12:deriv:inflection}{buigpunten}.

\emph{2.} In $a = 2$: $f(2) = 16 - 32 + 10 = -6$, $f'(2) = 32 - 48 = -16$,
\omterm{def:g12:deriv:tangent}{raaklijn} $y = -6 - 16(x-2) = -16x + 26$. De kloof is
\[
g(x) = f(x) + 16x - 26 = x^4 - 4x^3 + 16x - 16 .
\]
Omdat $g(2) = g'(2) = g''(2) = 0$, deelt $(x-2)^3$ de uitdrukking $g$; delen
geeft $g(x) = (x-2)^3(x+2)$. In de buurt van $x = 2$ is $x + 2 > 0$, dus
heeft $g$ het teken van $(x-2)^3$: negatief vóór $2$, positief erna. De
kromme kruist de \omterm{def:g12:deriv:tangent}{raaklijn}: $2$ is inderdaad een \omterm{def:g12:deriv:inflection}{buigpunt}.
\end{solution}

\begin{solution}{exo:g12:deriv:7}
De \omterm{def:g11:func:function}{functie} $f(x) = \sqrt{x}$ voldoet aan $f''(x) = -\frac{1}{4}x^{-3/2} < 0$
op $\intoo{0}{+\infty}$: ze is \omterm{def:g12:deriv:convex}{concaaf}, zodat haar \omterm{def:g10:functions:graph}{grafiek} \emph{onder} elke
\omterm{def:g12:deriv:tangent}{raaklijn} ligt. De \omterm{def:g12:deriv:tangent}{raaklijn} in $a = 1$ is
\[
y = f(1) + f'(1)(x - 1) = 1 + \tfrac12 (x-1) = \frac{x+1}{2},
\]
waaruit $\sqrt{x} \leq \frac{x+1}{2}$ voor alle $x > 0$. Gelijkheid betekent
dat de kromme de \omterm{def:g12:deriv:tangent}{raaklijn} raakt, en dat gebeurt alleen in het raakpunt
$x = 1$. (Gelijkwaardig: $\left(\sqrt x - 1\right)^2 \geq 0$.)
\end{solution}

\begin{solution}{exo:g12:deriv:8}
Volgens \cref{thm:g12:deriv:convexity} ligt de \omterm{def:g10:functions:graph}{grafiek} van een
differentieerbare convexe \omterm{def:g11:func:function}{functie} boven elk van haar \omterm{def:g12:deriv:tangent}{raaklijnen}. De \omterm{def:g12:deriv:tangent}{raaklijn}
in $a$ is horizontaal ($f'(a) = 0$), met \omterm{def:g10:algebra:equation}{vergelijking} $y = f(a)$. Bijgevolg is
$f(x) \geq f(a)$ voor alle $x \in \R$: $f(a)$ is het globale \omterm{def:g10:functions:extrema}{minimum}.
\end{solution}

\begin{solution}{exo:g12:deriv:9}
\emph{Vaste omtrek.} Een rechthoek met zijden $x$ en $\frac{p}{2} - x$
($0 < x < \frac p2$) heeft oppervlakte $S(x) = x\left(\frac p2 - x\right)$.
Dan verdwijnt $S'(x) = \frac p2 - 2x$ in $x = \frac p4$, positief ervóór en
negatief erna: het \omterm{def:g10:functions:extrema}{maximum} ligt in $x = \frac{p}{4}$, waar beide zijden
$\frac p4$ meten --- een vierkant.

\emph{Vaste oppervlakte.} Zijden $x$ en $\frac Ax$ geven omtrek
$P(x) = 2\left(x + \frac Ax\right)$, met
$P'(x) = 2\left(1 - \frac{A}{x^2}\right)$, negatief voor $x < \sqrt A$ en
positief erna: \omterm{def:g10:functions:extrema}{minimum} in $x = \sqrt{A}$, opnieuw een vierkant.

\emph{Verband.} De twee uitspraken zijn duaal. Stel dat een niet-vierkante
rechthoek bij vaste oppervlakte $A$ de omtrek minimaal maakte; het vierkant
met dezelfde omtrek zou volgens de eerste uitspraak oppervlakte $> A$ hebben,
en het gelijkvormig verkleinen tot oppervlakte $A$ zou zijn omtrek strikt
verkleinen --- in tegenspraak met de minimaliteit. Elke uitspraak houdt de
andere dus in.
\end{solution}

\begin{solution}{pb:g12:deriv:1}
\textbf{1.} $30x\left(3x^2 + 1\right)^4$; \quad
$\dfrac{x}{\sqrt{x^2 + 9}}$; \quad
$-\dfrac{2x}{\left(x^2 + 1\right)^2}$.

\textbf{2.} $y(4) = 5$, \omterm{def:g10:reffunc:affine}{richtingscoëfficiënt} $\frac45$: \omterm{def:g12:deriv:tangent}{raaklijn}
$y = 5 + \frac45(x - 4)$, dus $y = \frac45 x + \frac95$.

\textbf{3.} $f'(x) = \frac{(x^2 + 1) - x \cdot 2x}{(x^2+1)^2}
= \frac{1 - x^2}{(x^2 + 1)^2}$: negatief, positief, negatief over $-1$ en
$1$ heen: \omterm{def:g10:functions:extrema}{minimum} $f(-1) = -\frac12$, \omterm{def:g10:functions:extrema}{maximum} $f(1) = \frac12$, met
\omterm{def:g12:limcont:asymptote}{horizontale asymptoot} $0$ aan beide kanten.

\textbf{4.} $g'(x) = 3x^2 - 6x = 3x(x - 2)$: \omterm{def:g12:seq:monotonic}{stijgend}, \omterm{def:g10:functions:variations}{dalend}, \omterm{def:g12:seq:monotonic}{stijgend} over
$0$ en $2$ heen; \omterm{def:g12:deriv:extremum}{lokaal maximum} $g(0) = 4$, lokaal \omterm{def:g10:functions:extrema}{minimum} $g(2) = 0$.
$g''(x) = 6x - 6$: \omterm{def:g12:deriv:convex}{concaaf} vóór $1$, \omterm{def:g12:deriv:convex}{convex} erna: \omterm{def:g12:deriv:inflection}{buigpunt} in $(1, 2)$.

\textbf{5.} \omterm{def:g12:deriv:tangent}{Raaklijn} in $1$: \omterm{def:g10:reffunc:affine}{richtingscoëfficiënt} $g'(1) = -3$, dus
$y = 2 - 3(x - 1) = -3x + 5$. \omterm{def:g11:seq:arithmetic}{Verschil}:
$x^3 - 3x^2 + 4 - (-3x + 5) = x^3 - 3x^2 + 3x - 1 = (x - 1)^3$, dat in $1$
van teken verandert: de \omterm{def:g12:deriv:tangent}{raaklijn} \emph{kruist} de kromme --- het kenmerkende
gedrag in een \omterm{def:g12:deriv:inflection}{buigpunt}, waar de kromme van kant wisselt.

\textbf{6.} $T(x) = \dfrac{\sqrt{x^2 + 900}}{5} +
\dfrac{\sqrt{(40 - x)^2 + 400}}{2}$. Het water invóór $0$ of voorbij $40$
verlengt \emph{beide} stukken: zinloos.

\textbf{7.} $T'(x) = \dfrac{x}{5\sqrt{x^2 + 900}} -
\dfrac{40 - x}{2\sqrt{(40 - x)^2 + 400}}$.

\textbf{8.} In de loopdriehoek is de zijde langs de waterlijn $x$ en de
schuine zijde $\sqrt{x^2 + 900}$: de \omterm{def:g11:trigo:cossin}{sinus} van de hoek met de loodlijn is
precies hun quotiënt (overstaande zijde op schuine zijde); net zo is
$\sin\theta_2 = \frac{40 - x}{\sqrt{(40-x)^2 + 400}}$. Dan luidt
$T'(x) = 0$ als $\frac{\sin\theta_1}{5} = \frac{\sin\theta_2}{2}$ --- de wet
van Snellius, met de snelheden $5$ en $2$.

\textbf{9.} $T(24) \approx 20.5$ s; $T(40) = 10 + 10 = 20$ s;
$T(0) = 6 + \sqrt{2000}/2 \approx 28.4$ s. De rechte lijn verliest zelfs van
``helemaal langs het strand lopen en dan recht naar buiten zwemmen'' --- en
allebei verliezen ze van de gebroken weg.

\textbf{10.} \omterm{thm:g12:limcont:ivt}{Bisectie} op $T'$ (negatief in $24$, positief in $40$)
\omterm{def:g12:seq:limit}{convergeert} naar $x^* \approx 34$ m, met $T(x^*) \approx 19.5$ s: ongeveer
één seconde sneller dan de rechte lijn --- het \omterm{def:g11:seq:arithmetic}{verschil} tussen een redding en
een drama.

\textbf{11.} Elke deeltijd is een convexe \omterm{def:g11:func:function}{functie} van $x$ (een tak van de
vorm $\sqrt{\text{kwadratische uitdrukking}}$), dus is $T$ \omterm{def:g12:deriv:convex}{convex}, en een
kritiek punt van een convexe differentieerbare \omterm{def:g11:func:function}{functie} is haar globale
\omterm{def:g10:functions:extrema}{minimum} (\cref{exo:g12:deriv:8}): $x^*$ is niet zomaar stationair, hij is
onverslaanbaar.

\textbf{12.} Licht is trager in water; volgens het principe van Fermat buigt
de snelste weg van lucht naar water aan het oppervlak precies volgens de wet
van Snellius --- zodat stralen van een ondergedompeld potlood gebroken bij het
oog aankomen, en de hersenen, die rechte lijnen doortrekken, het potlood aan
de waterlijn gebroken zien.

\textbf{13.} $(x^2)'' = 2 > 0$: \omterm{def:g12:deriv:convex}{convex}. \omterm{def:g12:deriv:convex}{Convexiteit} in het \omterm{prop:g10:coordgeom:midpoint}{midden}:
$f\left(\frac{a+b}{2}\right) \leq \frac{f(a) + f(b)}{2}$, en dat is de
bewering. Met de hand:
$\frac{a^2 + b^2}{2} - \left(\frac{a+b}{2}\right)^2 = \frac{(a - b)^2}{4}
\geq 0$.

\textbf{14.} De vierkantswortel (\omterm{def:g12:seq:monotonic}{stijgend}) nemen in vraag 13 geeft
$\frac{a + b}{2} \leq \sqrt{\frac{a^2 + b^2}{2}}$. De volledige ketting op
$2$ en $8$: harmonisch $\frac{2 \cdot 16}{10} = 3.2$; \omterm{def:g12:seq:arith-geom}{meetkundig}
$\sqrt{16} = 4$; \omterm{def:g12:seq:arith-geom}{rekenkundig} $5$; kwadratisch $\sqrt{34} \approx 5.83$: elk
\omterm{def:g11:stat:mean}{gemiddelde} buigt voor het volgende.

\textbf{15.} \omterm{def:g12:deriv:tangent}{Raaklijn} aan $\frac1x$ in $1$: waarde $1$, \omterm{def:g10:reffunc:affine}{richtingscoëfficiënt}
$-1$: $y = 2 - x$. Convexe krommen liggen boven hun \omterm{def:g12:deriv:tangent}{raaklijnen}
(\cref{thm:g12:deriv:convexity}): $\frac1x \geq 2 - x$ voor alle $x > 0$, met
gelijkheid precies in het raakpunt $x = 1$.

\textbf{16.} In het \omterm{def:g12:deriv:inflection}{buigpunt} bereikt het \emph{dagelijkse} aantal gevallen
(de \omterm{def:g12:deriv:derivative}{afgeleide}) haar top: de groei houdt op te versnellen en begint te
vertragen --- het eerste wiskundige signaal dat de golf keert, lang voordat de
aantallen zelf dalen. Ervóór is $f'' > 0$ (elke dag erger dan de vorige);
erna is $f'' < 0$ (nog groeiend, maar trager).

\textbf{17.} $S'(r) = 4\pi r - \frac{2V}{r^2} = 0$ geeft
$r^3 = \frac{V}{2\pi} = \frac{330}{2\pi}$: $r \approx 3.74$ cm, en
$h = \frac{330}{\pi r^2} \approx 7.49$ cm.

\textbf{18.} In het optimum is $h = \frac{V}{\pi r^2} =
\frac{2\pi r^3}{\pi r^2} = 2r$: de hoogte is gelijk aan de diameter, wat het
volume ook is.

\textbf{19.} $S''(r) = 4\pi + \frac{4V}{r^3} > 0$: \omterm{def:g12:deriv:convex}{convex}, dus is de kritieke
straal het globale \omterm{def:g10:functions:extrema}{minimum}. Echte blikjes staan hoger omwille van de grip,
het stapelen, de zichtbaarheid in het rek en het dikkere metaal boven en
onder --- optimaliseren minimaliseert altijd precies wat je opschreef, niet
wat je bedoelde.

\textbf{20.} De redder: modelleer $T(x)$, leid af, Snellius, \omterm{def:g12:deriv:convex}{convexiteit}, één
seconde gewonnen. Het blikje: modelleer $S(r)$, leid af, $h = 2r$,
\omterm{def:g12:deriv:convex}{convexiteit}, metaal gespaard. De balk (\cref{pb:g11:deriv:1}): modelleer
$S(w)$, leid af, $d = w\sqrt2$, tabel, stijfheid gewonnen. En volgens het
principe van Fermat doorloopt het licht zelf die pijplijn aan elk oppervlak
dat het tegenkomt --- de natuur was de eerste optimaliseerder.
\end{solution}
